\documentclass[11pt]{amsart}
\usepackage[T1]{fontenc}
\usepackage{lmodern,amsmath,amssymb,amsthm,microtype}
\usepackage[margin=1.15in]{geometry}
\usepackage[colorlinks=true,linkcolor=blue,citecolor=blue,urlcolor=blue]{hyperref}
% Repository locator for the public (non-anonymous) build of
% non_mf_groups_exist.tex.  Kept in a separate file so that the anonymous
% build has an anonymous source archive and not merely an anonymous PDF:
% exclude this file from a double-blind submission and the manuscript still
% compiles, printing "supplied with this submission as supplementary
% material" in place of the link.
%
% Loaded only when \anonymousfalse is set in the preamble.
\newcommand{\coderepository}{https://github.com/SauersML/group-approximation}
\newcommand{\coderepositoryname}{github.com/SauersML/group-approximation}

\hypersetup{pdftitle={Isomorphisms between the simple Kazhdan groups of minimal subshifts},pdfauthor={\paperauthorname},pdfkeywords={simple groups, property (T), minimal subshifts, crossed products, strong orbit equivalence, elementary groups, Fredholm index}}
\newtheorem{theorem}{Theorem}[section]
\newtheorem{proposition}[theorem]{Proposition}
\newtheorem{lemma}[theorem]{Lemma}
\newtheorem{corollary}[theorem]{Corollary}
\theoremstyle{remark}
\newtheorem{remark}[theorem]{Remark}
\newcommand{\F}{\mathbb F}
\newcommand{\Z}{\mathbb Z}
\newcommand{\LC}{\operatorname{LC}}
\newcommand{\EL}{\operatorname{EL}}
\newcommand{\GL}{\operatorname{GL}}
\newcommand{\SL}{\operatorname{SL}}
\newcommand{\Fix}{\operatorname{Fix}}
\newcommand{\ind}{\operatorname{ind}}
\newcommand{\inn}{\operatorname{inn}}
\newcommand{\Out}{\operatorname{Out}}
\newcommand{\Aut}{\operatorname{Aut}}
\newcommand{\Hom}{\operatorname{Hom}}
\newcommand{\End}{\operatorname{End}}
\newcommand{\supp}{\operatorname{supp}}
\newcommand{\doi}[1]{\href{https://doi.org/#1}{\nolinkurl{doi:#1}}}
\newcommand{\arxiv}[1]{\href{https://arxiv.org/abs/#1}{\nolinkurl{arXiv:#1}}}
\raggedbottom
\emergencystretch=1.5em

% DRAFT, 2026-09-14. Transcribed by main from the reviewed Cairn chain
% (sk-orthogonality-o parts 1-2 with the merged fixes of sk-verify-25a/25b;
% sk-free-10 O1-O3 with sk-verify-18 part6; sk-free-6 Theorem D with
% sk-verify-15/16; rigid supports with sk-verify-3; Prop F with sk-verify-19;
% DP with sk-verify-16; K_0 with sk-verify-3). Citations marked TO VERIFY in
% the bibliography were not checked at source by the present author.

\begin{document}
\title{Isomorphisms between the simple Kazhdan groups of minimal subshifts}
\author{\paperauthorname}
\date{}
\subjclass[2020]{Primary 20E32; Secondary 16E20, 19B14, 20C20, 37B10}
\keywords{simple groups, property (T), minimal subshifts, crossed products, strong orbit equivalence, elementary groups}

\begin{abstract}
For an infinite minimal subshift $X$ let $R_X=\LC(X,\F_2)\rtimes\Z$ and
$G_X=\EL_3(R_X)$, the infinite simple Kazhdan LEF group of the companion
note. We show that every isomorphism $G_X\to G_Y$ is standard: it is
conjugation by an element of $\GL_3(R_Y)$ composed with a ring isomorphism
$R_X\to R_Y$ applied entrywise, possibly after the adjoint-inverse
automorphism. So $G_X\cong G_Y$ if and only if $R_X\cong R_Y$, and then
$X$ and $Y$ are strongly orbit equivalent. Along the way we show that
$\GL_3(R_X)=G_X\rtimes\Z$ through a Fredholm index, that every element of
$G_X$ is a product of at most $78$ elementary matrices, and that
$\Out(G_X)$ contains $(\Z\times\Aut(X,T))/\langle(-3,T)\rangle$.
\end{abstract}
\maketitle

\section{Introduction}\label{sec:intro}

Let $X\subseteq A^{\Z}$ be an infinite minimal subshift with shift $T$, let
$\LC(X,\F_2)$ be the ring of locally constant functions $X\to\F_2$, and let
\[
  R_X=\LC(X,\F_2)\rtimes_T\Z,\qquad G_X=\EL_3(R_X),
\]
where $R_X$ consists of the finite sums $\sum_jf_ju^j$ with
$ufu^{-1}=f\circ T^{-1}$, and $\EL_3$ is the group generated by the
elementary matrices $e_{ij}(r)=I_3+rE_{ij}$. The companion
note~\cite{Note} shows that $G_X$ is an infinite, finitely generated,
simple group with Kazhdan's property~\textup{(T)} that is a limit of finite
simple groups, and asks whether $G_X\cong G_Y$ forces $X$ and $Y$ to be flip
conjugate, or at least strongly orbit equivalent in the sense of Giordano,
Putnam and Skau~\cite{GPS}. We answer the second question.

Write $\iota$ for the anti-automorphism $\sum_jf_ju^j\mapsto\sum_ju^{-j}f_j$
of $R_X$, and for $g\in\GL_3(R_X)$ write $g^*$ for the transpose of the
matrix $\iota(g)$ obtained by applying $\iota$ to the entries. Then
$\gamma_X(g)=(g^*)^{-1}$ is an automorphism of $\GL_3(R_X)$ that preserves
$G_X$; we call it the adjoint-inverse automorphism. A ring isomorphism
$\varphi\colon R_X\to R_Y$ induces the isomorphism $M_3(\varphi)$ of matrix
rings, hence of $G_X$ onto $G_Y$, and conjugation by any $k\in\GL_3(R_Y)$
preserves $G_Y$ (Theorem~\ref{thm:index}). We call an isomorphism
$\alpha\colon G_X\to G_Y$ \emph{standard} if
$\alpha=\inn(k)\circ M_3(\varphi)|_{G_X}$ or
$\alpha=\inn(k)\circ M_3(\varphi)\circ\gamma_X|_{G_X}$ for such $k$ and
$\varphi$.

\begin{theorem}\label{thm:main}
Let $X$ and $Y$ be infinite minimal subshifts. Every isomorphism
$G_X\to G_Y$ is standard.
\end{theorem}

\begin{corollary}\label{cor:soe}
$G_X\cong G_Y$ if and only if $R_X\cong R_Y$ as rings, and then $(X,T)$ and
$(Y,T')$ are strongly orbit equivalent.
\end{corollary}

The proof has three parts. Section~\ref{sec:index} describes $\GL_3(R_X)$
through a Fredholm index: the elementary group is exactly the kernel of the
index, $\GL_3(R_X)=G_X\rtimes\langle\operatorname{diag}(u,1,1)\rangle$, and
every element of $G_X$ is a product of at most $78$ elementary matrices.
Section~\ref{sec:frob} looks at the copies of the Frobenius group
$F_{21}\le\GL_3(\F_2)$ that are constant on a clopen set $V$ and locates
their images under $\alpha$ by counting their double centralizers, which
have order~$168$ on both sides. This forces the images to act on
$R_Y^3$ through a single natural block, and the modular representation
theory of $\GL_3(\F_2)$ then shows that the blocks of disjoint clopen sets
are orthogonal. Section~\ref{sec:standard} turns these blocks into
idempotents of $R_Y$, using that $K_0(R_Y)$ is torsion-free and that every
idempotent of $M_3(R_Y)$ is similar to a diagonal one with clopen entries,
and recovers $\varphi$ from the root subgroups. Section~\ref{sec:k} collects
the $K$-theory of $R_X$ and the passage to strong orbit equivalence.

Flip conjugacy stays open: strong orbit equivalence classes contain systems
of different entropy, so no argument through $K_0$ alone can give more. For
topological full groups the converse is known~\cite{GPS99,BezuglyiMedynets},
and there the argument recovers the dynamics from a diagonal subalgebra
that isomorphisms preserve. Here an isomorphism need not preserve the
diagonal: conjugation by $\operatorname{diag}(u,1,1)$ is an automorphism of
$G_X$ that is not inner (Proposition~\ref{prop:out}) and does not preserve
$\GL_3(\LC(X,\F_2))$.

\section{The index and the elementary group}\label{sec:index}

Throughout, $(X,T)$ is a minimal homeomorphism of an infinite Cantor set,
$R=R_X$, $G=G_X$, and $e_U$ is the indicator of a clopen set $U$. We write
$d=\operatorname{diag}(u,1,1)$.

\subsection*{The orbit representation}
Fix $x\in X$ and let $V_x=\bigoplus_{t\in\Z}\F_2^3\delta_t$, on which
$f\in\LC(X,\F_2)$ acts by $f\delta_t=f(T^tx)\delta_t$ and $u$ by
$u\delta_t=\delta_{t+1}$; then $ufu^{-1}\delta_t=f(T^{t-1}x)\delta_t
=(f\circ T^{-1})\delta_t$, so this is a representation of $M_3(R)$ on
$V_x$. It is faithful, since the orbit of $x$ is dense, and every
$a\in M_3(R)$ acts by a banded matrix: its entry $a_x(m,n)\in M_3(\F_2)$
vanishes for $|m-n|>\deg a$, where $\deg a$ is the largest $|j|$ with
$u^j$ occurring in an entry of $a$, and it depends only on $T^mx$ through
finitely many clopen sets. For a half-line $J=(-\infty,c)$ or $J=[c,\infty)$
let $P_J$ be the coordinate projection onto $V_J=\bigoplus_{t\in J}\F_2^3\delta_t$.

\begin{lemma}\label{lem:fredholm}
Let $h\in\GL_3(R)$, let $w$ bound $\deg h$ and $\deg h^{-1}$, and let
$Q=P_Jh_xP_J$ on $V_J$ for $J=[c,\infty)$. Then $\ker Q$ is spanned by
vectors supported in $[c,c+2w)$, $Q$ maps onto every vector supported in
$[c+2w,\infty)$, and $Q$ is Fredholm. The index
$\ind_{x,c}(h)=\dim\ker Q-\dim\operatorname{coker}Q$ is a homomorphism
$\GL_3(R)\to\Z$, it vanishes on every $e_{ij}(r)$, and $\ind(d)=-1$.
\end{lemma}

\begin{proof}
If $Q\xi=0$, then $h_x\xi$ is supported outside $J$ and within $w$ of
$\supp\xi\subseteq J$, so in $[c-w,c)$, and $\xi=h_x^{-1}h_x\xi$ is
supported in $[c-2w,c+w)\cap J\subseteq[c,c+2w)$. If $\eta$ is supported in
$[c+2w,\infty)$, then $\xi=h_x^{-1}\eta$ is supported in $J$ and
$Q\xi=P_J\eta=\eta$. So kernel and cokernel are finite-dimensional. For
$a,b\in\GL_3(R)$, $P_Ja_xb_xP_J-(P_Ja_xP_J)(P_Jb_xP_J)=P_Ja_x(1-P_J)b_xP_J$
has finite rank, and the index of a Fredholm operator is additive and
invariant under finite-rank perturbations, so $\ind_{x,c}$ is a
homomorphism. For $a=e_{ij}(r)=I+N$ with $N^2=0$,
$Q^2=P_J+P_JNP_JNP_J=P_J-P_JN(1-P_J)NP_J$ is the identity up to finite
rank, so $2\ind(Q)=0$. Finally $d$ acts on the first coordinate as the
unilateral shift $\delta_t\mapsto\delta_{t+1}$ of $V_J$, which is injective
with cokernel spanned by $\delta_c$, and as the identity on the other two.
\end{proof}

\begin{lemma}\label{lem:constant}
$\ind_{x,c}(h)$ depends neither on $x$ nor on $c$.
\end{lemma}

\begin{proof}
Compressing a Fredholm operator to a subspace of finite codimension does
not change its index, as the compression differs from the operator
restricted to that subspace by finite rank; so $\ind_{x,c}=\ind_{x,c+1}$.
The basis change $\delta_t\mapsto\delta_{t-1}$ identifies $V_{Tx}$ with
$V_x$ and moves the cut $c$ to $c+1$, so $\ind_{Tx,c}=\ind_{x,c+1}$. By
Lemma~\ref{lem:fredholm} the kernel and a complement of the image of $Q$
lie in the window $[c,c+2w)$, and the index is a linear-algebra quantity
of the entries of $h_x$ and $h_x^{-1}$ on $[c-2w,c+4w)$, which depend only
on $T^tx$ for $t$ in a bounded window through finitely many clopen sets.
So $x\mapsto\ind_{x,c}(h)$ is locally constant and $T$-invariant, and
minimality makes it constant.
\end{proof}

We write $\ind(h)$ for this common value. For $g\in\GL_3(R)$ the two
half-line compressions at a cut differ from $g_x$ by the finite-rank corners
$P_Jg_x(1-P_J)$ and $(1-P_J)g_xP_J$, so their indices add to
$\ind(g_x)=0$: the compression to $(-\infty,c)$ has index $-\ind(g)$.

\begin{theorem}\label{thm:bounded}
Every $h\in\GL_3(R)$ with $\ind(h)=0$ is a product of at most $78$
elementary matrices, in a sequence of positions fixed in advance.
\end{theorem}

The proof cuts the orbit into long Kakutani--Rokhlin towers and compresses
$h$ to each tower. The vanishing index lets kernel and cokernel be matched
at each end of a tower separately, giving a block-diagonal factor with
banded inverse; the remaining factor is the identity away from the cuts and
is block diagonal for the towers based at midpoints. Both factors live in
finite products of matrix algebras over $\F_2$, where block Gauss
reduction, Whitehead diagonals and Thompson's commutator theorem give $39$
elementary matrices each. We first record the four tools.

\begin{lemma}[finite block Gauss reduction]\label{lem:T1}
Let $S=\prod_aM_{h_a}(\F_2)$ be a finite product with every $h_a\ge3$.
Every element of $\GL_3(S)$ is a product of at most $39$ elementary
matrices over $S$, in a fixed pattern of positions.
\end{lemma}

\begin{proof}
$S$ is semisimple Artinian, so it has stable rank one. Let
$A\in\GL_3(S)$ with first column $(a_1,a_2,a_3)$, which is left unimodular.
Stable rank one gives $t,t'\in S$ with $p=a_1+ta_3+t'a_2$ a unit, so
$e_{12}(t')e_{13}(t)A$ has the unit pivot $p$; two more left elementary
matrices clear the rest of the first column and two right ones the rest of
the first row, leaving $\operatorname{diag}(p,A')$ with $A'\in\GL_2(S)$. One
left, one left and one right elementary matrix reduce $A'$ to a diagonal.
So $A$ is a product of $9$ elementary matrices and
$\operatorname{diag}(p,q,r)=\operatorname{diag}(p,p^{-1},1)
\operatorname{diag}(1,pq,(pq)^{-1})\operatorname{diag}(1,1,pqr)$.
In characteristic $2$ each $\operatorname{diag}(c,c^{-1},1)$ is a product of
$6$ elementary matrices:
$e_{12}(c)e_{21}(c^{-1})e_{12}(c)$ is the matrix with $c$ and $c^{-1}$ off
the diagonal, and $e_{12}(1)e_{21}(1)e_{12}(1)$ is the transposition.
The unit $v=pqr$ lies in $\GL_1(S)=\prod_a\SL_{h_a}(\F_2)$, and by
Thompson's theorem~\cite{Thompson61} every element of $\SL_n(K)$ with
$n\ge3$ is a commutator, so $v=[x,y]$ componentwise with
$x,y\in\GL_1(S)$. Then
$\operatorname{diag}(1,1,[x,y])=\operatorname{diag}(1,x,x^{-1})
\operatorname{diag}(1,y,y^{-1})\operatorname{diag}(1,(yx)^{-1},yx)$ is a
product of $18$ elementary matrices. The total is $9+12+18=39$.
\end{proof}

\begin{lemma}[end correction]\label{lem:T3}
Let $h\in\GL_3(R)$ with $\ind(h)=0$, let $w$ bound $\deg h^{\pm1}$, put
$W=2w$, let $J=(-\infty,c)$ and $Q=P_Jh_xP_J$. There is a linear map $F$
on $V_J$, supported on the coordinates $[c-W,c)$ in rows and columns, such
that $Q+F$ is bijective and $(Q+F)^{-1}=P_Jh_x^{-1}P_J+D$ with $D$
supported in $[c-4W,c)^2$. The choice of $F$ depends only on the entries
of $h_x$ and $h_x^{-1}$ in the window $[c-3W,c+W)$. The same holds for
$J=[c,\infty)$ with the window $[c,c+4W)$.
\end{lemma}

\begin{proof}
By Lemma~\ref{lem:fredholm} applied to the left half-line, $K=\ker Q$ lies
in the span $L$ of the coordinates $[c-W,c)$, and $V_J=\operatorname{Im}Q+L$.
Choose a complement $L'$ of $K$ in $L$ and a subspace $M\subseteq L$ with
$V_J=\operatorname{Im}Q\oplus M$; then $\dim M=\dim\operatorname{coker}Q
=\dim K$, because the index is $0$. Let $F$ map $K$ isomorphically onto
$M$ and vanish on $L'$ and on the other coordinates. If $(Q+F)(k+\xi')=0$
with $k\in K$, then $Q\xi'=-Fk\in\operatorname{Im}Q\cap M=0$, so $k=0$ and
$\xi'\in\ker Q\cap(L'\oplus\text{far coordinates})=0$. So $Q+F$ is
injective, and being Fredholm of index $0$ it is bijective. Put
$B'=P_Jh_x^{-1}P_J$. Then $(Q+F)B'=P_J-P_Jh_x(1-P_J)h_x^{-1}P_J+FB'
=P_J+E$ with $E$ supported in $[c-2W,c)^2$, and likewise $B'(Q+F)=P_J+E''$.
So $D=(Q+F)^{-1}-B'=-(Q+F)^{-1}E=-E''(Q+F)^{-1}$ has its columns and its
rows in the window, and the window of $(Q+F)^{-1}E$ is at most $[c-4W,c)$.
All choices were made from the local entries.
\end{proof}

\begin{lemma}[two ends]\label{lem:T4}
In the setting of Lemma~\ref{lem:T3}, let $B=[c_0,c_1)$ with
$c_1-c_0\ge16W$, and let $F_{\mathrm{top}}$, $F_{\mathrm{bot}}$ be the
corrections of Lemma~\ref{lem:T3} for $(-\infty,c_1)$ and $[c_0,\infty)$.
Then $A_B=P_Bh_xP_B+F_{\mathrm{top}}+F_{\mathrm{bot}}$ is invertible on
$V_B$, with inverse $P_Bh_x^{-1}P_B+D_{\mathrm{top}}+D_{\mathrm{bot}}$.
\end{lemma}

\begin{proof}
Let $\tilde A$ be the displayed candidate and $n$ a coordinate in the upper
half of $B$. Then $D_{\mathrm{bot}}\delta_n=0$ and
$P_Bh_x^{-1}\delta_n=P_{(-\infty,c_1)}h_x^{-1}\delta_n$, so
$\xi=\tilde A\delta_n=(Q_{\mathrm{top}}+F_{\mathrm{top}})^{-1}\delta_n$,
which is supported within $4W$ of $n$ or of $c_1$, hence far from $c_0$.
So $P_Bh_x\xi=Q_{\mathrm{top}}\xi$ and $F_{\mathrm{bot}}\xi=0$, and
$A_B\xi=\delta_n$. The lower half is symmetric. So $A_B\tilde A=P_B$, and
$A_B$ is a square matrix.
\end{proof}

\begin{proof}[Proof of Theorem~\ref{thm:bounded}]
Let $h$, $w$, $W$ be as above and $L\ge\max(32W,3)$. Choose a
Kakutani--Rokhlin partition of $X$ into towers $T^iB_a$, $0\le i<h_a$, with
heights $h_a\ge L$, and refine the bases so that each base atom determines
the entries of $h_x$ and $h_x^{-1}$ along the whole occurrence of its tower
in every orbit, with a margin of $4W$ on both sides; finitely many atoms
suffice, because return times are bounded on a minimal system. Let
$A_{\mathcal P}\subseteq R$ be the span of the elements
$u^{i-i'}e_{T^{i'}B_a}$ with $0\le i,i'<h_a$; these multiply like matrix
units, so $A_{\mathcal P}\cong\prod_aM_{h_a}(\F_2)$, and $M_3(A_{\mathcal P})$
consists of the elements of $M_3(R)$ that act block diagonally on the
tower occurrences with entries depending only on the tower and the levels.

Let $H_1$ act on each tower occurrence $B$ as $A_B$ of Lemma~\ref{lem:T4}.
The corrections depend only on the local type, so $H_1\in M_3(A_{\mathcal P})$,
and its inverse of Lemma~\ref{lem:T4} lies there too. By Lemma~\ref{lem:T1},
$H_1$ is a product of at most $39$ elementary matrices over
$A_{\mathcal P}\subseteq R$. Now $Z=H_1^{-1}h-I$ acts as zero on every
$\delta_n$ farther than $7W$ from all cuts, and near a cut it is supported
within $7W$ of that cut. Let $\mathcal P'$ be the Kakutani--Rokhlin
partition over the midpoint levels $T^{\lfloor h_a/2\rfloor}B_a$, refined
in the same way; each of its occurrences contains exactly one cut of
$\mathcal P$, at distance at least $L/4$ from its ends. So
$H_2=I+Z$ and $H_2^{-1}=h^{-1}H_1$, which is localized in the same way,
lie in $M_3(A_{\mathcal P'})$, and $H_2$ is a product of at most $39$
elementary matrices. Then $h=H_1H_2$.
\end{proof}

\begin{theorem}\label{thm:index}
$\ker(\ind)=\EL_3(R)=G$. Consequently $\GL_3(R)=G\rtimes\langle d\rangle$
with $\GL_3(R)/G\cong\Z$, every element of finite order of $\GL_3(R)$ lies
in $G$, and $G=[\GL_3(R),\GL_3(R)]$.
\end{theorem}

\begin{proof}
Lemma~\ref{lem:fredholm} gives $G\subseteq\ker(\ind)$ and
Theorem~\ref{thm:bounded} the converse. For $k\in\GL_3(R)$,
$kd^{\ind(k)}$ has index $0$, so it lies in $G$, and
$\langle d\rangle\cap G=1$ since $\ind(d^m)=-m$. A homomorphism to $\Z$
kills torsion, and $[\GL_3(R),\GL_3(R)]\subseteq\ker(\ind)=G=[G,G]$, as
$G$ is simple and nonabelian~\cite{Note}.
\end{proof}

\begin{proposition}\label{prop:normal}
The normal subgroups of $\GL_3(R)$ are $1$ and $\ind^{-1}(m\Z)$ for
$m\ge0$. The center of $\GL_3(R)$ is trivial, and $Z(R)=\F_2$.
\end{proposition}

\begin{proof}
Let $1\ne g$ lie in a subgroup $K$ normalized by $G$. The simplicity proof
of~\cite{Note} applies word for word to $g\in\GL_3(R)$: it uses only the
entries of $g$ and $g^{-1}$ to produce a small clopen $V$ and an
$h=e_{ij}(e_V)$ not commuting with $g$; the commutator $[g,h]=g(hg^{-1}h^{-1})$
lies in $K$, since $h\in G$ normalizes $K$, and it lies in a finite simple
group $H_V\cong\GL_d(\F_2)\le G$ with $d\ge3$, so $H_V\subseteq K$; the set
$J$ of $r$ with $e_{pq}(r)\in K$ for all $p\ne q$ is a two-sided ideal
containing $e_V$, since $e_{pq}(sr)=[e_{pl}(s),e_{lq}(r)]$ and $e_{pl}(s)\in G$
normalizes $K$; and minimality gives $J=R$. So $K\supseteq G$, and the
normal subgroups containing $G$ correspond to the subgroups of
$\GL_3(R)/G\cong\Z$. For the center of $R$: if $r=\sum_nc_nu^n$ commutes
with every $f\in\LC(X,\F_2)$, then $c_n(f\circ T^{-n}-f)=0$, and for
$n\ne0$ aperiodicity gives a clopen set separating $x$ from $T^nx$ near any
point of $\supp c_n$, so $c_n=0$; then $c_0\circ T^{-1}=c_0$ makes $c_0$
constant. A matrix commuting with every $e_{ij}(1)$ is scalar, and a
scalar commuting with every $e_{ij}(r)$ has entry in $Z(R)^\times=1$.
\end{proof}

\begin{proposition}\label{prop:out}
Conjugation induces an injective homomorphism $\GL_3(R)/G\cong\Z\to\Out(G)$.
For $\phi\in\Aut(X,T)$, the automorphism $\alpha_\phi(\sum f_nu^n)=\sum(f_n\circ\phi^{-1})u^n$
of $R$ induces an automorphism of $G$, and
$(m,\phi)\mapsto[\inn(d^m)\circ\alpha_\phi]$ is a homomorphism
$\Z\times\Aut(X,T)\to\Out(G)$ with kernel $\langle(-3,T)\rangle$.
\end{proposition}

\begin{proof}
By Proposition~\ref{prop:normal} the centralizer of $G$ in $\GL_3(R)$ is
trivial, so $k\mapsto\inn(k)|_G$ is injective, and $\inn(k)$ is inner in
$G$ exactly when $k\in G$. Since $\alpha_\phi$ fixes $d$, the two kinds of
automorphisms commute. If $\inn(d^m)\alpha_\phi$ is inner, then
$\alpha_\phi=\inn(k)$ on $G$ with $k=d^{-m}g$, $g\in G$; as $\alpha_\phi$
fixes every $e_{ij}(1)$ and $e_{ij}(u)$, $k=cI_3$ with $c\in R^\times$
commuting with $u$. The centralizer of $u$ in $R$ is
$\{\sum c_nu^n:c_n\circ T=c_n\}=\F_2[u^{\pm1}]$ by minimality, whose units
are the $u^n$, so $\alpha_\phi(f)=u^nfu^{-n}=f\circ T^{-n}$ and $\phi=T^n$.
Additivity of the index with $\ind(u^nI_3)=-3n$ and $\ind(d^{-m}g)=m$ gives
$m=-3n$. Conversely $\inn(d^{-3n})\alpha_{T^n}=\inn(d^{-3n}u^nI_3)$ and
$d^{-3n}u^nI_3$ has index $0$, so lies in $G$.
\end{proof}

\section{Frobenius subgroups and their double centralizers}\label{sec:frob}

\subsection*{Supports}
For a clopen $U\subseteq X$ let $E_U=e_UI_3$, let
$G_U=G\cap(I_3+M_3(e_URe_U))$ be the elements supported in $U$, let
$L=\GL_3(\LC(X,\F_2))$, which is the group of locally constant maps
$X\to\GL_3(\F_2)$, and for a subgroup $Q\le\GL_3(\F_2)$ let
$Q_U=\{I_3+e_U(q-I_3):q\in Q\}$ be its copy constant on $U$ and trivial
off $U$. We use the copies of $Q=\GL_3(\F_2)$, of order $168$, and of its
Frobenius subgroup $F_{21}=\langle c,s\rangle\cong C_7\rtimes C_3$, where
$c$ is a Singer cycle with characteristic polynomial $x^3+x+1$.

\begin{lemma}[rigid supports]\label{lem:supports}
Let $U\subseteq X$ be clopen and nonempty.
\begin{enumerate}
\item[\textup{(a)}] The corner ring $e_URe_U$ is simple with center
$\F_2e_U$, and $\LC(U,\F_2)$ is maximal commutative in it.
\item[\textup{(b)}] $C_G(G_U)=G_{X\setminus U}$ and $C_G(L_{X\setminus U})=G_U$;
in particular $C_G(L)=1$.
\item[\textup{(c)}] $\operatorname{span}\{g-I:g\in G_U\}=M_3(e_URe_U)$.
\end{enumerate}
\end{lemma}

\begin{proof}
(a) $R$ is simple~\cite[Corollary~4.6]{ClarkEdie}, a corner of a simple
ring is simple, and $e_U$ is full, so $Z(e_URe_U)=e_UZ(R)=\F_2e_U$ by
Proposition~\ref{prop:normal}. If $c=\sum_jf_ju^j\in e_URe_U$ commutes
with every $e_V$, $V\subseteq U$ clopen, then
$f_j(e_V-e_{T^jV})=0$, and aperiodicity gives $f_j=0$ for $j\ne0$, so
$c\in\LC(U,\F_2)$.
(c) $e_{ij}(r)\in G_U$ for $r\in e_URe_U$ and $i\ne j$, the identity
$(g-I)(h-I)=(gh-I)-(g-I)-(h-I)$ makes the span a ring, and
$(rE_{ij})(sE_{ji})=rsE_{ii}$.
(b) If $h$ centralizes $G_U$, it commutes with $M_3(e_URe_U)\ni E_U$ by
(c), so it is block diagonal for $E_U$, and its $U$-corner is central in
$M_3(e_URe_U)$, hence $cE_U$ with $c\in\F_2$; invertibility gives $c=1$
and $h\in G_{X\setminus U}$. Conversely, elements supported on disjoint
clopen sets commute. For the second statement,
$\operatorname{span}(\GL_3(\F_2)-I)$ is a nonzero ideal of $M_3(\F_2)$, so
$\operatorname{span}(L_{X\setminus U}-I)=M_3(e_{X\setminus U}\LC(X,\F_2))$,
and a centralizer $h$ of $L_{X\setminus U}$ is block diagonal with
$X\setminus U$ corner $cI_3$, $c$ commuting with $\LC(X\setminus U,\F_2)$;
by (a), $c\in\LC(X\setminus U,\F_2)$ is invertible, so $c=e_{X\setminus U}$
and $h\in G_U$. With $U=\varnothing$, $C_G(L)=1$.
\end{proof}

\subsection*{Representations of $\GL_3(\F_2)$ over $\F_2$}
The simple $\F_2Q$-modules are the trivial module $1$, the natural module
$3$, its dual $3^*$, and the Steinberg module $8$, which is projective
\cite{ModularAtlas}. The group $F_{21}$ has odd order, so $\F_2F_{21}$ is
semisimple; its simple modules are $1$, the module $2$ inflated from $C_3$,
with endomorphism ring $\F_4$, and the restrictions $3$ and $3^*$ of the
natural module and its dual. Restrictions of semisimple modules are
determined by Brauer characters. On the classes $1,3A,7A,7B$ of $F_{21}$,
with $\alpha=(-1+\sqrt{-7})/2$, so that $\alpha+\bar\alpha=-1$ and
$\alpha\bar\alpha=2$, the Brauer characters are
$1=(1,1,1,1)$, $2=(2,-1,2,2)$, $3=(3,0,\alpha,\bar\alpha)$,
$3^*=(3,0,\bar\alpha,\alpha)$, and the Steinberg module of $Q$ has
character $(8,-1,1,1)$ on the classes $1A,3A,7A,7B$ of $Q$. We use:
\begin{enumerate}
\item[\textup{(F1)}] $8|_{F_{21}}=2\oplus3\oplus3^*$, and $8$ has no
nonzero $F_{21}$-fixed vector.
\item[\textup{(F2)}] $3\otimes3^*=1\oplus8$ and
$\operatorname{Ext}^1_{\F_2Q}(3,3)=0$. Indeed $3\otimes3^*=\End(3)$ splits
as $\F_2I\oplus\mathfrak{sl}_3$ because $\operatorname{tr}I=1$, and
$\mathfrak{sl}_3$ has the Brauer character of $8$; then
$\operatorname{Ext}^1(3,3)=H^1(Q,1)\oplus H^1(Q,8)=0$, since $Q$ is
perfect and $8$ is projective.
\item[\textup{(F3)}] Restricted to the diagonal $F_{21}$,
$3\otimes3=3\oplus3^*\oplus3^*$, since $\alpha^2=\bar\alpha-1=\alpha+2\bar\alpha$.
\item[\textup{(F4)}] $\F_2[F_{21}]$ maps onto $M_3(\F_2)$ through the
natural module, which is absolutely irreducible.
\end{enumerate}

\begin{lemma}\label{lem:S3}
$3\otimes3$ is not a semisimple $\F_2Q$-module, and neither is
$3^*\otimes3^*$.
\end{lemma}

\begin{proof}
The symmetric tensors $\operatorname{Sym}\subseteq3\otimes3$ form a
submodule of dimension $6$ containing
$\Lambda=\operatorname{span}\{v\otimes w+w\otimes v\}\cong\Lambda^2(3)\cong3^*$,
and $v\mapsto v\otimes v$ induces an isomorphism
$\operatorname{Sym}/\Lambda\cong3$, additive in characteristic $2$. If
$\operatorname{Sym}=\Lambda\oplus L$ for a submodule $L$, let
$s\colon3\to L$ be the equivariant section and $q(v)=s(v)+v\otimes v\in\Lambda$.
Then $q$ is equivariant, and its polarization is
$q(v+w)+q(v)+q(w)=v\otimes w+w\otimes v$. The stabilizer of $e_1$, the
matrices with first column $e_1$, fixes no nonzero vector of $3^*$: a fixed
functional $(f_1,f_2,f_3)$ must satisfy $f_1b=0$ for all $b\in\F_2^2$ and
$(f_2,f_3)A=(f_2,f_3)$ for all $A\in\GL_2(\F_2)$. So $q(e_1)=0$, by
transitivity $q=0$, and the polarization vanishes, a contradiction. The
adjoint-inverse automorphism exchanges $3$ and $3^*$.
\end{proof}

\subsection*{The double centralizer on the $X$ side}

\begin{lemma}\label{lem:Z}
Let $S$ be a simple ring with center $F$, and let $f\in S$ be an idempotent
with $f\ne0,1$. If $x\in S$ commutes with $fS(1-f)$ and $(1-f)Sf$, then
$x\in F$.
\end{lemma}

\begin{proof}
Write $x$ in Peirce blocks. From $xn=nx$ for $n\in fS(1-f)$,
$x_{(1-f)f}fS(1-f)=0$, and fullness of $1-f$ gives $x_{(1-f)f}=0$;
symmetrically $x_{f(1-f)}=0$. Then $x_{ff}n=nx_{(1-f)(1-f)}$ and
$x_{(1-f)(1-f)}m=mx_{ff}$ for $m\in(1-f)Sf$, so $x_{ff}$ commutes with
$fS(1-f)Sf=fSf$ and is central in $fSf$; as $f$ is full,
$Z(fSf)=fF$, so $x_{ff}=\lambda f$, and likewise
$x_{(1-f)(1-f)}=\lambda'(1-f)$. Then $\lambda n=\lambda'n$ for all
$n\in fS(1-f)\ne0$, so $\lambda=\lambda'$.
\end{proof}

\begin{lemma}\label{lem:X}
For every nonempty clopen $V$, $C_G(C_G((F_{21})_V))=Q_V$.
\end{lemma}

\begin{proof}
Write $M_X=R^3$ for the natural right $R$-module. The group $(F_{21})_V$
acts on it with isotypic parts $e_VM_X$, of type $3$, and $e_{X\setminus V}M_X$,
trivial, and $\End_{F_{21}}(3)=\F_2$; so every $g\in C_G((F_{21})_V)$ is
$rI_3$ on $e_VM_X$ with $r\in(e_VRe_V)^\times$ and arbitrary on
$e_{X\setminus V}M_X$. Hence $Q_V$ commutes with $C_G((F_{21})_V)$.
Conversely let $x$ be in the double centralizer. It commutes with
$G_{X\setminus V}\subseteq C_G((F_{21})_V)$, so $x\in G_V$ by
Lemma~\ref{lem:supports}(b). For orthogonal nonzero idempotents $f+f'=e_V$
of $R_V=e_VRe_V$ and $n\in fR_Vf'$, the element
$(e_V+n)I_3+E_{X\setminus V}$ has order $2$, so it lies in $G$ by
Theorem~\ref{thm:index}, and it commutes with $(F_{21})_V$, as $n$ is a
scalar on $V$. So every entry of $x-I$ commutes with $fR_Vf'$ and
$f'R_Vf$, and Lemma~\ref{lem:Z} for the simple ring $R_V$ with center
$\F_2e_V$ (Lemma~\ref{lem:supports}(a)) puts it in $\F_2e_V$. So $x\in Q_V$.
\end{proof}

\subsection*{The double centralizer on the $Y$ side}
From now on $\alpha\colon G_X\to G_Y$ is an isomorphism, $M=R_Y^3$ is the
natural right $R_Y$-module, on which $\GL_3(R_Y)$ acts by $R_Y$-linear
maps, and $K_V=\alpha((F_{21})_V)$. Since $|F_{21}|$ is odd, $M$ is the
direct sum of its $K_V$-isotypic parts $M_t=T_tM$, $t\in\{1,2,3,3^*\}$,
where the $T_t$ are the images of the central idempotents of
$\F_2F_{21}$, which lie in the span of $K_V$; each $M_t$ is an
$R_Y$-submodule, $M_3\cong3\otimes_{\F_2}P_3$ with
$P_3=\Hom_{F_{21}}(3,M)$ a right $R_Y$-module, likewise $M_{3^*}$, and
$M_2\cong2\otimes_{\F_4}P_2$. We put
\[
  S_V=I+\sum_{k\in K_V}k,
\]
the idempotent projecting onto $M_2\oplus M_3\oplus M_{3^*}=[M,K_V]$ along
$M_1=\Fix(K_V)$; on the $X$ side the same expression gives $E_V$, because
a Singer cycle has no nonzero fixed vector.

\begin{theorem}\label{thm:C}
For every nonempty clopen $V\subseteq X$ there is exactly one
$t_V\in\{3,3^*\}$ with $M_{t_V}\ne0$, and $M_2=0$. Moreover $\alpha(Q_V)$
acts trivially on $\Fix(K_V)=M_1$, and on $M_{t_V}=t_V\otimes P$ it acts as
$\psi_V(q)\otimes\mathrm{id}_P$ for an isomorphism
$\psi_V\colon Q\to\GL(t_V)$ that extends the representation $t_V$ of
$F_{21}$.
\end{theorem}

\begin{proof}
By Lemma~\ref{lem:X} and since $\alpha$ is an isomorphism,
$C_{G_Y}(C_{G_Y}(K_V))=\alpha(Q_V)$ has order $168$. Every
$x\in C_{G_Y}(K_V)$ commutes with each $T_t$, and on
$M_3\cong3\otimes P_3$ it lies in
$\End_{F_{21}\otimes R_Y^{\mathrm{op}}}(3\otimes P_3)=\mathrm{id}\otimes\End_{R_Y}(P_3)$;
likewise on $M_{3^*}$, and on $M_2$ it is $\F_4$-linear. So for
$a\in\GL_3(\F_2)$ the map $g_3(a)$ that is $a\otimes\mathrm{id}$ on $M_3$
and the identity on the other parts commutes with $C_{G_Y}(K_V)$. It is
$R_Y$-linear, so it lies in $\GL_3(R_Y)$, and it has finite order, so it
lies in $G_Y$ by Theorem~\ref{thm:index} for $Y$. Hence
$g_3(\GL_3(\F_2))\subseteq\alpha(Q_V)$; the same holds for $g_{3^*}$ and
for multiplication by $\omega\in\F_4^\times$ on $M_2$. The element
$\alpha(c_V)\ne I$ acts trivially on $M_1\oplus M_2$, so
$M_3\oplus M_{3^*}\ne0$. If both were nonzero, or if $M_2\ne0$, then
$\alpha(Q_V)$ would contain $\GL_3(\F_2)^2$ or $\GL_3(\F_2)\times C_3$, of
order more than $168$. So exactly one $t_V$ occurs and $M_2=0$, and
$g_{t_V}(\GL_3(\F_2))=\alpha(Q_V)$ by the order count. That is the
displayed form, and $\psi_V$ restricts to $t_V$ on $F_{21}$ because
$K_V\subseteq\alpha(Q_V)$ acts as $t_V\otimes\mathrm{id}$.
\end{proof}

\begin{corollary}\label{cor:O}
$S_VS_W=0$ for disjoint clopen $V,W$.
\end{corollary}

\begin{proof}
$Q_V$ and $Q_W$ commute, so $\alpha(Q_W)$ commutes with $\alpha(Q_V)$ and
with $K_V$; it therefore preserves $M_{t_V}=t_V\otimes P$ and acts there as
$\mathrm{id}\otimes\sigma$. So $N=S_VS_WM$, the part of $M_{t_V}$ on which
$K_W$ acts with type $t_W$, is a module $t_V\boxtimes t_W\otimes P_N$ for
$\alpha(Q_V)\times\alpha(Q_W)$, and the diagonal $\alpha(Q_{V\sqcup W})$ acts
on it as $(t_V\otimes t_W)\otimes\mathrm{id}$. By Theorem~\ref{thm:C} at
$V\sqcup W$, $M$ is a semisimple $\alpha(Q_{V\sqcup W})$-module with
constituents $1$ and $t_{V\sqcup W}$, and so is its submodule $N$. But
$3\otimes3$ and $3^*\otimes3^*$ are not semisimple (Lemma~\ref{lem:S3}),
and $3\otimes3^*=1\oplus8$ has the constituent $8$. So $P_N=0$.
\end{proof}

\begin{corollary}\label{cor:Oprime}
$S_X=I$, and $t_V=t_X$ for every nonempty clopen $V$.
\end{corollary}

\begin{proof}
Let $V,W$ be disjoint. By Corollary~\ref{cor:O}, $M=S_VM\oplus S_WM\oplus M'$
with $M'=(I-S_V-S_W)M$; $K_V$ acts trivially on the last two summands and
without fixed vectors on the first, and $K_W$ symmetrically, so the diagonal
$K_{V\sqcup W}$ has fixed space $M'$ and $S_{V\sqcup W}=S_V+S_W$. Put
$E=I-S_X$. Then $E\le I-S_V$ for every $V$, so $\alpha(Q_V)$ is trivial on
$EM\subseteq\Fix(K_V)$ (Theorem~\ref{thm:C}), and it preserves
$S_XM=S_VM\oplus S_{X\setminus V}M$, since it commutes with $K_{X\setminus V}$.
The group $L_X$ is generated by the $Q_V$, so $\alpha(L_X)$ is trivial on
$EM$ and preserves $S_XM$. For $n\in EM_3(R_Y)E$ with $n^2=0$, the element
$I+n$ commutes with $\alpha(L_X)$ and lies in $G_Y$ by
Theorem~\ref{thm:index}; but $C_{G_Y}(\alpha(L_X))=\alpha(C_{G_X}(L_X))=1$
by Lemma~\ref{lem:supports}(b), so $n=0$. If $E\ne0$, then by
Theorem~\ref{thm:DP} $E$ is similar to a diagonal idempotent with clopen
entries, so it splits into two nonzero orthogonal idempotents $f,f'$, and
$fM_3(R_Y)f'\ne0$ by simplicity; any $n$ in it has $n^2=0$ and lies in
$EM_3(R_Y)E$. So $E=0$.

For the orientation, $S_VM=(I-S_{X\setminus V})M=\Fix(K_{X\setminus V})$,
on which $\alpha(Q_{X\setminus V})$ is trivial, so $\alpha(Q_X)=\alpha(Q_V)\alpha(Q_{X\setminus V})$
acts there as $\alpha(Q_V)$, with $K_X$ of type $t_V$. By
Theorem~\ref{thm:C} at $X$, $K_X$ has type $t_X$ on all of $M$, so
$t_V=t_X$.
\end{proof}

\section{Standardness}\label{sec:standard}

Write $\rho_3(q)=q$ and $\rho_{3^*}(q)=(q^t)^{-1}$ for the natural
representation of $Q=\GL_3(\F_2)$ and its dual, and $\gamma(q)=(q^t)^{-1}$,
so that $\rho_{3^*}\circ\gamma=\rho_3$. The map $\psi_X$ of
Theorem~\ref{thm:C} is a faithful representation of $Q$ on a
$3$-dimensional space that restricts to $t_X$ on $F_{21}$; the only such
representations are $3$ and $3^*$, whose restrictions to $F_{21}$ are
inequivalent, so $\psi_X\cong\rho_{t_X}$, and the intertwiner commutes with
$F_{21}$, hence is scalar by (F4): $\psi_X=\rho_{t_X}$. If $t_X=3^*$, then
$(\alpha\circ\gamma_X)(q_X)=\alpha(\gamma(q)_X)=\rho_{3^*}(\gamma(q))\otimes\mathrm{id}
=q\otimes\mathrm{id}$, so the isomorphism $\alpha\circ\gamma_X$ has
orientation $3$, and Section~\ref{sec:frob} applies to it. Replacing
$\alpha$ by $\alpha\circ\gamma_X$ if necessary, we assume $t_X=3$ from now
on. By Corollary~\ref{cor:Oprime}, $\Fix(K_X)=0$, so Theorem~\ref{thm:C}
at $X$ says that $M\cong3\otimes_{\F_2}P$ with $P=\Hom_{F_{21}}(3,M)$ and
that $\alpha(q_X)$ acts as $q\otimes\mathrm{id}$ for every $q\in Q$.

\begin{proof}[Proof of Theorem~\ref{thm:main}]
\emph{Coordinates.} As right $R_Y$-modules, $P^3\cong M=R_Y^3$, so $P$ is
finitely generated projective and $3[P]=3[R_Y]$ in $K_0(R_Y)$, which is
torsion-free (Theorem~\ref{thm:K}); so $[P]=[R_Y]$. Since $P$ is a summand
of $R_Y^3$, it is the image of an idempotent of $M_3(R_Y)$, which by
Theorem~\ref{thm:DP} is similar to $\operatorname{diag}(e_{W_1},e_{W_2},e_{W_3})$
for clopen $W_i\subseteq Y$; so $P\cong\bigoplus_ie_{W_i}R_Y$ with
$\sum_i[\chi_{W_i}]=[\chi_Y]$ in $K^0(Y,T')=C(Y,\Z)/(1-T'_*)$. By
Matui's division lemma~\cite[Lemma~2.5]{Matui05}, $[\chi_U]\le[\chi_W]$
if and only if some $\gamma$ in the topological full group of $(Y,T')$
maps $U$ into $W$, and $[\chi_U]=[\chi_W]$ if and only if some $\gamma$
maps $U$ onto $W$. So there are $\gamma_2,\gamma_3$ with
$W_2'=\gamma_2W_2\subseteq Y\setminus W_1$ and
$\gamma_3W_3=Y\setminus(W_1\sqcup W_2')$, and since
$e_WR_Y\cong e_{\gamma W}R_Y$ through the unit of $R_Y$ implementing
$\gamma$, we get $P\cong R_Y$. So $M\cong\F_2^3\otimes R_Y=R_Y^3$ by an
$R_Y$-linear isomorphism intertwining $\alpha(Q_X)$ with the constant
$\GL_3(\F_2)$: there is $k\in\GL_3(R_Y)$ such that
$\alpha'=\inn(k^{-1})\circ\alpha$ satisfies $\alpha'(q_X)=q$ for every
$q\in\GL_3(\F_2)$. By Theorem~\ref{thm:index}, $\alpha'$ maps $G_X$ onto
$G_Y$.

\emph{Scalar supports.} Put $S'_V=k^{-1}S_Vk$. It commutes with
$\alpha'((F_{21})_V)$ and with $\alpha'((F_{21})_{X\setminus V})$, hence
with the constant $F_{21}$, whose span is $M_3(\F_2)$ (F4); the commutant
of $M_3(\F_2)$ in $M_3(R_Y)$ is $R_YI_3$, so $S'_V=b_VI_3$ for an idempotent
$b_V\in R_Y$. By Corollaries~\ref{cor:O} and~\ref{cor:Oprime},
$V\mapsto S_V$ is additive on disjoint unions with $S_X=I$; it is injective,
since $S_V\ne0$ for $V\ne\varnothing$, and it preserves meets, by
expanding $S_V$ and $S_W$ over $V\cap W$, $V\setminus W$ and $W\setminus V$.
So $V\mapsto b_V$ is an injective unital homomorphism of Boolean algebras
into the commuting idempotents of $R_Y$, and
$\mathfrak B=\operatorname{span}\{b_V\}\cong\LC(X,\F_2)$.

\emph{The diagonal.} By Theorem~\ref{thm:C}, $\alpha(Q_V)$ is trivial on
$(I-S_V)M$ and preserves $S_VM$, so
$\alpha'(q_V)-I\in b_VM_3(R_Y)b_V$, and likewise
$\alpha'(q_{X\setminus V})-I\in(1-b_V)M_3(R_Y)(1-b_V)$. Their product is
$\alpha'(q_X)=q$ and the cross term vanishes, so
$\alpha'(q_V)=I+b_V(q-I)$. As $L_X$ is generated by the $q_V$,
$\alpha'(L_X)=\GL_3(\mathfrak B)$, and for $i\ne j$,
$\alpha'(\{e_{ij}(f):f\in\LC(X,\F_2)\})=\{I+bE_{ij}:b\in\mathfrak B\}$.

\emph{Root subgroups.} Let $A_{ij}=\{e_{ij}(r):r\in R_X\}$. A matrix
$g\in\GL_3(R_X)$ commutes with $e_{ij}(f)$ for all $f\in\LC(X,\F_2)$ exactly
when $g_{pi}=0$ for $p\ne i$, $g_{jq}=0$ for $q\ne j$, and
$g_{ii}=g_{jj}$ commutes with $\LC(X,\F_2)$; by
Lemma~\ref{lem:supports}(a) this entry lies in $\LC(X,\F_2)$, and it is a
unit, so it equals $1$. The same conditions describe the centralizer of
$A_{ij}$, so $C_{G_X}(A_{ij}\cap L_X)=C_{G_X}(A_{ij})$, and the
double centralizer of $A_{ij}$ is $A_{ij}$ itself: intersecting the
conditions for $A_{12}$, $A_{13}$ and $A_{32}$, all of which lie in
$C(A_{12})$, leaves exactly the matrices $e_{12}(r)$. Applying $\alpha'$,
\[
  \alpha'(A_{ij})=C_{G_Y}C_{G_Y}(\{I+bE_{ij}:b\in\mathfrak B\})
  \subseteq C_{G_Y}C_{G_Y}(A_{ij}(R_Y))=A_{ij}(R_Y),
\]
using $Z(R_Y)=\F_2$ for the last equality. So
$\alpha'(e_{ij}(r))=e_{ij}(\varphi_{ij}(r))$ for additive injections
$\varphi_{ij}\colon R_X\to R_Y$ with images $J_{ij}\ni1$, since $\alpha'$
fixes $e_{ij}(1)$. From $[e_{ij}(r),e_{jk}(s)]=e_{ik}(rs)$ we get
$\varphi_{ik}(rs)=\varphi_{ij}(r)\varphi_{jk}(s)$, so $J_{ij}J_{jk}\subseteq J_{ik}$,
and with $1\in J_{jk}$ all $J_{ij}$ equal one subring $J$. As
$G_Y=\alpha'(G_X)$ is generated by the $e_{ij}(J)\subseteq M_3(J)$, every
$e_{12}(r)$ lies in $M_3(J)$, so $J=R_Y$.

\emph{Standardness.} Taking $s=1$ and $r=1$ in
$\varphi_{ik}(rs)=\varphi_{ij}(r)\varphi_{jk}(s)$ gives
$\varphi_{ik}=\varphi_{ij}=\varphi_{jk}$, so all $\varphi_{ij}$ equal one
bijection $\varphi\colon R_X\to R_Y$, which is additive and multiplicative:
a ring isomorphism. Then $\alpha'$ and $M_3(\varphi)$ agree on the
generators $e_{ij}(r)$ of $G_X$, so $\alpha=\inn(k)\circ M_3(\varphi)|_{G_X}$.
\end{proof}

\begin{proof}[Proof of Corollary~\ref{cor:soe}]
A ring isomorphism $R_X\to R_Y$ induces $G_X\cong G_Y$. Conversely, by
Theorem~\ref{thm:main} an isomorphism gives a ring isomorphism
$R_X\to R_Y$, or one $R_X\to R_Y$ composed with the adjoint-inverse
automorphism; in the second case $\iota$ still gives $R_X\cong R_Y$, as
$\iota$ is an anti-automorphism of $R_X$. A ring isomorphism $R_X\cong R_Y$
induces an isomorphism of unital ordered groups
$(K^0(X,T),K^0(X,T)^+,[1])\cong(K^0(Y,T'),K^0(Y,T')^+,[1])$
(Theorem~\ref{thm:K}), and Giordano, Putnam and Skau~\cite{GPS} showed that
this is equivalent to strong orbit equivalence.
\end{proof}

\section{$K$-theory of the crossed product}\label{sec:k}

Let $(X,T)$ be a minimal Cantor system, $A=\LC(X,\F_2)$ and
$R=A\rtimes_T\Z=A[u,u^{-1};\sigma]$ with $\sigma(f)=f\circ T^{-1}$. Write
$K^0(X,T)=C(X,\Z)/(1-T_*)C(X,\Z)$ with $T_*f=f\circ T^{-1}$, the positive
cone $\{[f]:f\ge0\}$ and the unit $[1]$; this is the dimension group of
Herman, Putnam and Skau~\cite{HPS}.

\begin{theorem}\label{thm:K}
The inclusion $A\subseteq R$ induces an isomorphism
$K_0(R)\cong K^0(X,T)$ of abelian groups with $[e_U]\mapsto[\chi_U]$, under
which the classes of idempotents form the positive cone. The group
$K^0(X,T)$ is torsion-free.
\end{theorem}

\begin{proof}
$A$ is the directed union of the finite products of fields $A_{\mathcal P}\cong\F_2^{\mathcal P}$
over the finite clopen partitions $\mathcal P$, and every $A[t_1,\dots,t_p]$
is a directed union of finite products of polynomial rings over $\F_2$ with
flat transition maps, so $A$ is regular supercoherent. Since $\sigma$ is a
unital automorphism, the $K$-theory fibration for twisted Laurent
polynomial rings of Ara, Brustenga and Corti\~nas~\cite[Theorem~3.6 and Lemma~7.2]{ABC}
gives an exact sequence
\[
  K_0(A)\xrightarrow{1-\sigma_*}K_0(A)\to K_0(R)\to K_{-1}(A)\xrightarrow{1-\sigma_*}K_{-1}(A).
\]
$K$-theory preserves filtered colimits, so $K_0(A)=C(X,\Z)$ with
$[e_U]\mapsto\chi_U$, and $K_{-1}(A)=0$ because $K_{-1}$ of a regular
Noetherian ring vanishes~\cite{Bass}. The image of $1-\sigma_*$ equals that
of $1-T_*$, which gives the group isomorphism.

Nonnegative functions are sums of clopen indicators, so their classes are
classes of idempotents. Conversely let $e\ne0$ be an idempotent of some
$M_n(R)$ and $\mu$ an ergodic invariant probability measure. Ara and
Claramunt~\cite{AraClaramunt} construct a faithful Sylvester matrix rank
function $\operatorname{rk}_\mu$ on $R$ with $\operatorname{rk}_\mu(e_U)=\mu(U)$;
it induces a state $s_\mu$ of $K_0(R)$ with $s_\mu[\chi_U]=\mu(U)$, hence
$s_\mu=\int\cdot\,d\mu$, and $\int[e]\,d\mu=\operatorname{rk}_\mu(e)>0$. By
Bauer's maximum principle the same holds for every invariant measure, and
since $K^0(X,T)$ is a simple dimension group whose states are the
integrals against invariant measures~\cite{HPS}, Effros's positivity
criterion~\cite{Effros} puts $[e]$ in the positive cone.

If $nf=g-g\circ T^{-1}$ with $n\ge1$, then $g$ modulo $n$ is $T$-invariant,
so constant by minimality; writing $g=c+nh$ gives $f=h-h\circ T^{-1}$. So
$K^0(X,T)$ is torsion-free.
\end{proof}

\begin{theorem}\label{thm:DP}
Every idempotent $e\in M_n(R)$ is similar, by a unit of $M_n(R)$, to a
diagonal idempotent $\operatorname{diag}(e_{W_1},\dots,e_{W_n})$ with
clopen sets $W_i\subseteq X$.
\end{theorem}

\begin{proof}
Represent $M_n(R)$ faithfully on $V=\bigoplus_{t\in\Z}\F_2^n\delta_t$ as in
Section~\ref{sec:index}, and let $w$ bound $\deg e$. Fix a cut $c$, and let
$V_\pm$ be the spans of the $\delta_t$ with $t<c$ and $t\ge c$. Put
$E=eV$, $E'=(1-e)V$, $E_\pm=E\cap V_\pm$ and $E'_\pm=E'\cap V_\pm$. For
$|t-c|\ge w$ the vector $\delta_t=e\delta_t+(1-e)\delta_t$ splits on one
side, so $V_\pm=E_\pm\oplus E'_\pm\oplus M_\pm$ with $M_\pm$ spanned by the
coordinates in $[c-w,c)$ and $[c,c+w)$; the sum
$\Sigma=E_-\oplus E'_-\oplus E_+\oplus E'_+$ is direct, because
$E\cap E'=0$ and $V_-\cap V_+=0$. Let $\tilde I$ and $\tilde I'$ be the
images of $E$ and $E'$ in $V/\Sigma\cong M:=M_-\oplus M_+$; they span, and
they are independent, since $\xi-\eta\in\Sigma$ with $\xi\in E$, $\eta\in E'$
forces $\xi\in\Sigma$. By a Laplace expansion of the invertible matrix of
the two projections along the coordinate basis $B$ of $M$, there is a
subset $T\subseteq B$ with $\operatorname{span}T\oplus\tilde I'=M$ and
$\operatorname{span}(B\setminus T)\oplus\tilde I=M$. Let $e_c$ be the
projection onto $E_-\oplus E_+\oplus\operatorname{span}T$ along
$E'_-\oplus E'_+\oplus\operatorname{span}(B\setminus T)$. It agrees with
$e$ on $\Sigma$, its range and kernel split along $V_-\oplus V_+$, and
$\Delta_c=e_c-e$ is supported in the window $[c-w,c+w)^2$ with entries
depending only on the local configuration. The intertwiner
$h_c=e_ce+(1-e_c)(1-e)=1+\Delta_c(2e-1)$ satisfies $h_ce=e_ch_c$ and is
invertible, with a window-local inverse: it maps $E$ onto the range of
$e_c$ and $E'$ onto its kernel, bijectively in both cases.

Choose a Kakutani--Rokhlin partition with towers of height more than
$4w+2$ and base atoms refined to fix the local types, and apply the
construction at every cut, with disjoint windows. Then $e''=e+\sum_c\Delta_c$
is an idempotent, since $e\Delta_c+\Delta_ce+\Delta_c^2=\Delta_c$ and the
windows are disjoint; it splits at every cut, so it lies in the tower
algebra $A_{\mathcal P}\cong\prod_aM_{h_a}(\F_2)$ tensored with $M_n$, and
$h=1+\sum_c\Delta_c(2e-1)$ is a unit with $e''=heh^{-1}$. In
$M_n(\prod_aM_{h_a}(\F_2))$ an idempotent of rank $k_a$ on the tower $a$ is
similar to a coordinate idempotent, which is the indicator of a union of
levels, so $e''$ is similar to $\operatorname{diag}(e_{W_1},\dots,e_{W_n})$.
\end{proof}

\section*{Origin and authorship}
Claude (Anthropic) found the construction and original proofs under the
author's direction on September~14, 2026, and wrote the manuscript. Proof
records are in the
\href{https://github.com/SauersML/group-approximation}{project repository}.

\begin{thebibliography}{99}

% TO VERIFY at source before submission: ABC (theorem numbers 3.6, 7.2 read
% from the e-print by the un-lane), AraClaramunt (e-print), Bass (locator),
% Effros (locator), HPS (locator), Matui05 (arXiv read; journal data),
% ModularAtlas (page), Thompson61 (through Urschel).

\bibitem{ABC}
P.~Ara, M.~Brustenga, and G.~Corti\~nas, \emph{$K$-theory of Leavitt path
algebras}, M\"unster J. Math. \textbf{2} (2009), 5--33. \arxiv{0903.0056}.

\bibitem{AraClaramunt}
P.~Ara and J.~Claramunt, \emph{Sylvester matrix rank functions on crossed
products}, Ergodic Theory Dynam. Systems \textbf{40} (2020), 2913--2946.
\arxiv{1902.06476}.

\bibitem{Bass}
H.~Bass, \emph{Algebraic $K$-theory}, W. A. Benjamin, New York, 1968.

\bibitem{BezuglyiMedynets}
S.~Bezuglyi and K.~Medynets,
\emph{Full groups, flip conjugacy, and orbit equivalence of Cantor minimal
systems}, Colloq. Math. \textbf{110} (2008), 409--429.
\doi{10.4064/cm110-2-6}.

\bibitem{ClarkEdie}
L.~O. Clark and C.~Edie-Michell, \emph{Uniqueness theorems for Steinberg
algebras}, Algebr. Represent. Theory \textbf{18} (2015), 907--916.
\doi{10.1007/s10468-015-9522-2}.

\bibitem{Effros}
E.~G. Effros, \emph{Dimensions and $C^*$-algebras}, CBMS Regional Conf.
Ser. in Math. \textbf{46}, Amer. Math. Soc., Providence, 1981.

\bibitem{GPS}
T.~Giordano, I.~F. Putnam, and C.~F. Skau,
\emph{Topological orbit equivalence and $C^*$-crossed products},
J. Reine Angew. Math. \textbf{469} (1995), 51--111.
\doi{10.1515/crll.1995.469.51}.

\bibitem{GPS99}
T.~Giordano, I.~F. Putnam, and C.~F. Skau,
\emph{Full groups of Cantor minimal systems},
Israel J. Math. \textbf{111} (1999), 285--320.
\doi{10.1007/BF02810689}.

\bibitem{HPS}
R.~H. Herman, I.~F. Putnam, and C.~F. Skau, \emph{Ordered Bratteli
diagrams, dimension groups and topological dynamics}, Internat. J. Math.
\textbf{3} (1992), 827--864. \doi{10.1142/S0129167X92000382}.

\bibitem{Matui05}
H.~Matui, \emph{Approximate conjugacy and full groups of Cantor minimal
systems}, Publ. Res. Inst. Math. Sci. \textbf{41} (2005), 695--722.
\arxiv{math/0404224}.

\bibitem{ModularAtlas}
C.~Jansen, K.~Lux, R.~Parker, and R.~Wilson, \emph{An atlas of Brauer
characters}, London Math. Soc. Monogr. (N.S.) \textbf{11}, Oxford Univ.
Press, New York, 1995.

\bibitem{Note}
\paperauthorname, \emph{Infinite simple Kazhdan groups that are limits of
finite simple groups}, manuscript (2026).

\bibitem{Thompson61}
R.~C. Thompson, \emph{Commutators in the special and general linear
groups}, Trans. Amer. Math. Soc. \textbf{101} (1961), 16--33.
\doi{10.1090/S0002-9947-1961-0130917-7}.

\end{thebibliography}
\end{document}
